\documentclass{template}

\usepackage{mathtools}
\pgfplotsset{compat=1.18}

\newcommand{\qedwhite}{\hfill \ensuremath{\Box}}


\begin{document}

\begin{center}
    \Large \textcolor{blue!70!black}{\textbf{Finite Field Extensions: Existence and Uniqueness}} \\
    \large \textbf{Mathematics Seminar at Kyiv School of Economics} \\
    \normalsize \textit{Dmytro Zakharov}
\end{center}

\begin{abstract}
    This small lecture is devoted to a quick glance at the structure of finite
    fields, primarly based on~\cite[Chapter 7]{ireland1990classical}. One
    obvious well-known example of the finite field is $\mathbb{Z}/p\mathbb{Z}$,
    but is this the only possible finite field? We first show the explicit
    construction of the finite field of order $p^n$ for prime $p$ and build the
    intuition of how elements of such fields look like. We then show that
    \textit{(i)} $p^n$ is the only possible such a finite field order,
    \textit{(ii)} always possible regardless of the chosen characteristic $p$,
    and \textit{(iii)} such construction is unique (up to an isomorphism).
\end{abstract}

\tableofcontents

\newpage

\section{Introduction}

In this section, we \textit{informally} introduce finite field extensions and
specify how to think of them. We will demonstrate how to build quadratic field
extensions and show how one constructs higher-degree field extensions.

\subsection{Basic Notions}

For the sake of completeness, we recall what is a field and introduce some basic
notation.

\begin{definition}[Finite Field]
    Recall that $K$ is a \textbf{field} if:
    \begin{itemize}
        \item $(K,+)$ and $(K\setminus\{0\},\times)$ form
        abelian groups. We denote the latter group simply by
        $K^{\times}$.
        \item Distributive law holds: $a(b+c)=ab+ac$ for arbitrary $a,b,c \in K$.
    \end{itemize}
    Field $K$ is \textbf{finite} if it contains the finite number of elements
    $q$. In such case, one commonly denotes such a field by $\mathbb{F}_q$ or
    even $\text{GF}(q)$ (latter is more common in Computer Science).
\end{definition}

The most basic example of such a field is the \textbf{prime field}
$\mathbb{Z}/p\mathbb{Z}$ or simply $\mathbb{F}_p$. For instance, for $p=5$, one
can perform operations over $\mathbb{F}_5$ as follows:
\begin{gather*}
    2+4 = 1, \quad 3 \cdot 4 = 2, \quad 3^{-1} = 2 \quad \textcolor{gray!75!black}{(\textit{all over $\mathbb{F}_5$})}
\end{gather*}

\textbf{Motivating Question.} Is this the only form of a finite field?

\subsection{Quadratic Field Extensions}
Of course the answer is negative. Let us give the example of the field with
$p^2$ elements.

\begin{example}
    Suppose it so happens that $p \equiv 3 \pmod{4}$. Then, the 
    following quotient ring
    \begin{equation*}
        \mathbb{F}_{p^2} \cong \mathbb{F}_p[x] / \langle x^2+1 \rangle
    \end{equation*}
    is a finite field with $p^2$ elements.
\end{example}

\textbf{Proof.} Recall that $K[x]/\langle f(x) \rangle$ is a field iff
$f(x) \in K[x]$ is irreducible over field $K$. Thus it
suffices to show that $x^2+1$ is irreducible over $\mathbb{F}_p$ for $p \equiv 3
\pmod{4}$. This reduces down to showing that $-1$ is not a quadratic residue in
$\mathbb{F}_p$. Accoring to Euler's criterion, one has
$\left(\frac{-1}{p}\right) \equiv (-1)^{(p-1)/2} \pmod{p}$. Since $p=4m+3$, we
obtain $\left(\frac{-1}{p}\right) = (-1)^{2m+1} = -1$, so $-1$ is a quadratic
non-residue. This proves that $\mathbb{F}_p[x] / \langle x^2+1 \rangle$ is a
field. Since such field consists of linear polynomials $\alpha+\beta x$ with
$\alpha,\beta \in \mathbb{F}_p$, we conclude that such finite field indeed
consists of $p^2$ elements. \qedwhite

\begin{example}
    In fact, $\mathbb{F}_{p^2}$ greatly resembles complex numbers $\mathbb{C}$
    (especially if we use $i$ to denote the variable $x$). For instance, letting
    $p=19$, operations in $\mathbb{F}_{361} \cong \mathbb{F}_{19}(i)$ look as
    follows:
    \begin{align*}
        \textbf{Addition:} && (1+10i) + (18+15i) = 19+25i = 6i \\
        \textbf{Multiplication:} && (5+6i)(6+7i) = 30+71i+42i^2 = 7+14i.
    \end{align*} 
\end{example}

\begin{remark}
    In fact, for other primes $p>2$ (such that $p \equiv 1 \pmod{4}$) one might
    use the similar construction: $\mathbb{F}_{p^2} \cong
    \mathbb{F}_p[x]/\langle x^2+\alpha \rangle$ for quadratic non-residue
    $\alpha$. It is well-known that for every $p>2$ we can find at least one
    quadratic non-residue (moreover, there are exactly $\frac{p-1}{2}$ of them).
\end{remark}

In fact, $\mathbb{F}_{p^2}$ is a field extension of $\mathbb{F}_p$. An extension
of a field $K$ is a field $L \supset K$. Note that this inclusion is a morphism
of fields, which particularly means that the identity $e$ of $K$ is also the
identity of $L$. Note also that $L$ is naturally a vector space over $K$ because
for any $\alpha,\beta \in K$ and $x,y \in L$, one can form linear combinations
$\alpha x+\beta y \in L$. One denotes by $[L:K]$ the dimension of $L$ viewed as
a $K$-vector space. A notation $L/K$ is often used for a field extension $K
\subset L$.

As an example, $\mathbb{F}_{p^2}$ constructed above is a field extension of
$\mathbb{F}_p$ since \textit{(i)} $\mathbb{F}_p \subset \mathbb{F}_{p^2}$ and
\textit{(ii)} $\mathbb{F}_{p^2}$ is a field. Therefore, it is a linear space
with scalars from $\mathbb{F}_p$. Specifically, the basis is given by $\{1,i\}$.
In Proposition~\ref{proposition:alphas}, we will formalize this intuition more
systematically and formally.

\subsection{Explicit Construction for $\mathbb{F}_{p^n}$}

Now we notice that extending the previous construction to the $n$-degree
extension should be straightforward. Indeed, fix some polynomial $f \in
\mathbb{F}_p[x]$ that is irreducible over $\mathbb{F}_p$ and such that $\deg
f=n$. Then, the finite field extension $\mathbb{F}_{p^n}$ is defined as:
\begin{equation*}
    \mathbb{F}_{p^n} \cong \mathbb{F}_p[x] / \langle f(x) \rangle,
\end{equation*}
so elements of $\mathbb{F}_{p^n}$ are polynomials $\sum_{j=0}^{n-1}\alpha_jx^j$
with $\alpha_j \in \mathbb{F}_p$ and arithmetic is done ``modulo $f(x)$''.

The core three questions here are: 
\begin{enumerate}
    \item[\textbf{Q1}] Can we find any other finite fields of order other than $p^n$?
    \item[\textbf{Q2}] Does such $f$ always exist?
    \item[\textbf{Q3}] Can we construct another field $K$ of order $p^n$ such
    that $K \not\cong \mathbb{F}_p[x] / \langle f(x) \rangle$?
\end{enumerate}

All questions are in fact crucial. For instance, in the informal discussion
above, we were quite sketchy with the notation $\mathbb{F}_{p^n}$. Suppose we
had two finite fields, say, $K_0$ and $K_1$, of order $p^n$ such that $K_0
\not\cong K_1$. In such case, it would be unclear what we denote by
$\mathbb{F}_{p^n}$: should it be the field $K_0$, the field $K_1$, or some class
of such fields?

Since we give the positive answer to \textbf{Q2} and negative to \textbf{Q3},
such an issue with the definition does not arise. Specifically, we prove both
facts in \Cref{section:proving}. But first, we need more fundamental
mathematical preliminaries on finite fields, which we specify in
\Cref{section:basic-properties}.

% \subsection{Some Practical Applications*}
% Field Extensions play a key role in cryptographic protocols. In this section, we
% specify some of the interesting use cases where field extensions appear. The
% reader should feel free to skip this section and navigate to
% \Cref{section:proving} for the main Mathematical details.

% \subsubsection{Increasing Domain Size}

% Many cryptographic protocols operate with finite fields $K$: for instance, it is
% very convenient to encode large set of data into the single polynomial in
% $K[x]$. However, if the field $K$ is too-small, the security of the protocol
% reduces. In fact, the probability of some adversary breaking the cryptographic
% protocol is typically bounded above as $N/\#K$ where $N$ is the size of the
% computational problem. For that reason, if $K=\mathbb{F}_p$, 

\section{Basic Properties of Finite Fields}\label{section:basic-properties}

In this section, we analyze finite fields more systematically and rigorously.
Further always assume $K=\mathbb{F}_q$ is a finite field with order $q$. We
shall first prove that $q=p^n$ for some prime $p$ and natural $n$. Additionally,
we analyze the structure of multiplicative and additive groups of $K$.

\begin{lemma}\label{lemma:galois-polynomial}
    For any finite field $K$ of order $q$ one has:
    \[x^q-x = \prod_{\alpha \in K}(x-\alpha)\]
\end{lemma}

\textbf{Proof.} Notice that every $\alpha \in K$ is a root of $x^q-x$. Indeed,
if $\alpha=0$, the result is obvious. Otherwise,
$\alpha^q-\alpha=\alpha(\alpha^{q-1}-1)$ and since $K^{\times}$ is of order
$q-1$, this equates to zero as well. Thus, $\prod_{\alpha \in K}(x-\alpha)$
divides $x^q-x$. Since both polynomials are of degree $q$ and the leading 
coefficient of both expressions is $1$, we obtain the equality, as required. \qedwhite

This result is fundamental for our further discussion. In particular, we provide
the following two important corollaries.

\begin{proposition}\label{proposition:galois-polynomial-corollary}
    The following holds:
    \begin{enumerate}
        \item[\textit{(i)}] Let $L/K$ be the field extension. An element $\alpha \in K$ iff
    $\alpha^q=\alpha$.
        \item[\textit{(ii)}] For $f \in K[x]$, let $f(x) \mid (x^q-x)$ and $d := \deg
        f$. Then, $f$ has $d$ distinct roots.
    \end{enumerate}
\end{proposition}

\textbf{Proof.} 

\textit{Part (i).} $\alpha^q=\alpha$ iff it is a root of $x^q-x$. Since all
roots of $x^q-x$ are exactly elements of $K$ from
Lemma~\ref{lemma:galois-polynomial}, the result follows.

\textit{Part (ii).} Suppose $f(x)g(x) = x^q-x$ with $\deg g = q-d$. If 
$f$ had less than $d$ distinct roots, then $f(x)g(x)$ would have fewer 
than $q$ roots, which is not the case since $x^q-x$ has $q$ roots. \qedwhite

\subsection{Properties of Multiplicative Subgroup}

Now we shall discuss the properties of $K^{\times}$. We first prove the 
following fact.
\begin{theorem}\label{theorem:fp-cyclic}
    The multiplicative group of a finite field $K^{\times}$ is \textit{cyclic}.
\end{theorem}

\textbf{Proof.} Fix $d \mid (q-1)$ and let $K_d$ denote the set of elements of
$K^{\times}$ of order $d$. Suppose $K_d \neq \emptyset$, so there exists some
$\beta \in K_d$. Clearly, $\langle \beta \rangle \subseteq \{\alpha \in
K^{\times}: \alpha^d=1\}$ (since for any $\beta^i \in \langle \beta \rangle$,
one has $(\beta^i)^d = (\beta^d)^i=1$). But notice that $\langle \beta \rangle$
has a cardinality $d$, so $\langle \beta \rangle = \{\alpha \in K^{\times}:
\alpha^d=1\}$. Thus, $K_d$ is a set of generators of $\langle \beta \rangle$ of
order $d$, so $\# K_d = \varphi(d)$, where $\varphi$ is the Euler totient
function.

So we have proven that either $\# K_d = 0$ or $\# K_d = \varphi(d)$. So we shall
prove that the latter is true for any $d \mid (q-1)$. Notice that we can count
the number of elements in $K^{\times}$ as $\#K^{\times} = \sum_{d \mid (q-1)}\#
K_d$. On the other hand, $\sum_{d \mid (q-1)}\#K_d \leq \sum_{d \mid
(q-1)}\varphi(d) = q-1$. Thus, $\#K_d=\varphi(d)$. In particular, by letting
$d:=q-1$, one has $\#K_{q-1}=\varphi(q-1)>0$, so the generator of $K^{\times}$
is any $\omega \in K_{q-1}$. \qedwhite

\begin{definition}
    The generators of $K^{\times}$ are called \textbf{primitive elements}.
\end{definition}

Two important consequences of this fact are the following.
\begin{proposition}
    If $d \mid (q-1)$, then there is a cyclic subgroup $G \leq K^{\times}$ of
    order $d$. If $\omega$ is the primitive element of $K$, the generator of $G$
    is $\omega^{(q-1)/d}$.
\end{proposition}

\begin{proposition}
    Let $\alpha \in K^{\times}$. Then equation $x^n=\alpha$ has solutions iff
    $\alpha^{(q-1)/d}=1$, where $d = \gcd(n,q-1)$. If there are solutions, then
    there are exactly $d$ solutions.
\end{proposition}

\textbf{Proof.} Since $K^{\times}$ is cyclic, let $\langle \omega \rangle =
K^{\times}$. Then, $\alpha = \omega^a$ for some $a$. Since $x=0$ is clearly not
a solution, then $x \in K^{\times}$, so let $x=\omega^y$. Thus, we have an
equation $\omega^{ny} = \omega^a$ with respect to $y$. This clearly reduces down
to equation $ny \equiv a \pmod{q-1}$. It is then well-known that this equation
is solvable iff $d = \gcd(n,q-1) \mid a$ and the number of solutions is $d$. But
$d \mid a$ iff $\alpha^{(q-1)/d}=1$.\qedwhite

We have investigated the structure of $K^{\times}$. Let us now turn our
attention of the additive subgroup of $K$ which would give us information on the
possible values of $\#K$.

\subsection{Properties of Additive Subgroup}

\begin{lemma}\label{lemma:identity-subfield}
    Let $K$ be a finite field. Denote by $\langle e \rangle_+$ the set,
    additively generated by the identity $e$ of $K^{\times}$. Then, $\langle e
    \rangle_+$ is a subfield of $K$, isomorphic to $\mathbb{Z}/p\mathbb{Z}$ for
    some prime number $p$.
\end{lemma}

\textbf{Proof.} Define morphism $\phi: \mathbb{Z} \to K, n \mapsto ne$. Note
that $\langle e \rangle_+ = \text{im}(\phi)$. We also notice that $\phi$ is
clearly a ring homomorphism. Therefore, $\text{im}(\phi)$ forms a finite subring
of $K$, so in particular, it is an integral domain. Recall that
$\mathbb{Z}/\text{ker}(\phi) \cong \text{im}(\phi)$ due to first isomorphism
theorem, so in particular $\mathbb{Z}/\text{ker}(\phi)$ is an integral domain.
Therefore, $\text{ker}(\phi)$ is a prime ideal, so $\text{ker}(\phi) =
p\mathbb{Z}$ for some prime $p$. Thus we conclude $\langle e \rangle_+ =
\text{im}(\phi) \cong \mathbb{Z}/p\mathbb{Z}$.\qedwhite

\begin{definition}
    Let $K$ be a finite field and $e$ be the identity of $K^{\times}$. The (prime)
    number $p$ is called a \textbf{characteristic} of $K$ (and denoted as
    $\text{char}(K)$) if $\langle e \rangle_+ \cong \mathbb{Z}/p\mathbb{Z}$.
\end{definition}

We shall finally prove the theorem on the number of elements in $K$.
\begin{proposition}
    Let $K$ be a finite field. Then, $\#K = p^n$ where $p$ is a characteristic
    of $K$.
\end{proposition}

\textbf{Proof.} From Lemma~\ref{lemma:identity-subfield}, we know that
$\mathbb{F}_p$ is a subfield of $K$ for some prime $p$. Notice that $K$ is a
$\mathbb{F}_p$-vector space. Let $[K:\mathbb{F}_p]=n$. Then, there exists a
basis, say, $\beta_1,\dots,\beta_n \in K$, such that any $\beta \in K$ can be
written as $\beta = \sum_{i=1}^n \alpha_i\beta_i$ with $\alpha_i \in
\mathbb{F}_p$. There are clealy $p^n$ elements in $K$.\qedwhite

\begin{proposition}[Freshman's Dream]\label{proposition:freshman-dream}
    Suppose $K$ has a characteristic of $p$. Then,
    $(\alpha+\beta)^{p^d}=\alpha^{p^d} + \beta^{p^d}$ for any finite field
    elements $\alpha,\beta \in K$ and any positive integer $d$.
\end{proposition}

\textbf{Proof.} We prove this fact by induction. For $d=1$, notice
\begin{equation*}
    (\alpha+\beta)^p = \sum_{j=0}^p \begin{pmatrix}
        p \\ j
    \end{pmatrix}\alpha^j\beta^{p-j}.
\end{equation*}

Notice that for each $j \neq 0$ and $j \neq p$, one has $p \mid \left(\begin{smallmatrix}
p \\ j \end{smallmatrix}\right)$. Thus, all terms except for $j=0$ and $j=p$ vanish, 
so we conclude $(\alpha+\beta)^p = \alpha^p + \beta^p$.

Now suppose this formula holds for $d$, that is
$(\alpha+\beta)^{p^d}=\alpha^{p^d} + \beta^{p^d}$. Raise both sides by the power
of $p$. In such case, left-hand side is exactly $(\alpha+\beta)^{p^{d+1}}$. The
right-hand side, after applying the result from the induction base, one gets
$\alpha^{p^{d+1}}+\beta^{p^{d+1}}$.
\qedwhite

\section{Uniqueness and Existence}\label{section:proving}

\subsection{Preliminaries}

\begin{proposition}\label{proposition:ell-m}
    Let $K$ be a field. Then, $(x^{\ell}-1) \mid (x^m-1)$ in $K[x]$ iff
    $\ell \mid m$.
\end{proposition}

\textbf{Proof.} We first show that $\ell \mid m$ implies $(x^{\ell}-1) \mid
(x^m-1)$. If $\ell \mid m$, then $m=n\ell$ for some integer $n$. Then,
\begin{equation*}
    \frac{x^m-1}{x^{\ell}-1} = \frac{x^{n\ell}-1}{x^{\ell}-1} = \sum_{j=0}^{n-1}x^{j\ell} \in K[x]
\end{equation*}

Now we show the opposite direction. Suppose $m=n\ell + r$ with $0 \leq r <
\ell$. Then we have
\begin{equation*}
    \frac{x^m-1}{x^{\ell}-1} = \frac{x^{n\ell+r}-1}{x^{\ell}-1} = \frac{x^{n\ell}x^r-1}{x^{\ell}-1} = x^r \cdot \frac{x^{n\ell} - 1}{x^{\ell}-1} + \frac{x^r-1}{x^{\ell}-1}
\end{equation*}

As we have already shown, $(x^{n\ell}-1)/(x^{\ell}-1) \in K[x]$, so in order for the
left-hand side to be a polynomial, $(x^r-1)/(x^{\ell}-1)$ must be a polynomial.
Since degree of $x^r-1$ is $r$, which is less than the degree in the denominator
(which is $\ell$), the only case when $(x^{\ell}-1) \mid (x^r-1)$ is $r=0$.
\qedwhite

We shall finally show that for $d$, there is an in irreducible 
polynomial in $\mathbb{F}_p[x]$ of degree $d$. We first provide 
and show the more general fact. 
\begin{proposition}
    Suppose $K$ is a field and $f \in K[x]$ is irreducible. There exists some
    other field $L \supset K$ and an element $\alpha \in L$ such that
    $f(\alpha)=0$.
\end{proposition}

\textbf{Proof.} Recall that $K[x]$ is a principal ideal domain. Therefore,
$\langle f(x) \rangle$ is a maximal ideal and therefore $K[x]/\langle f(x)
\rangle$ is a field. Let $K' := K[x]/\langle f(x) \rangle$ and let $\iota: K[x]
\to K'$ be the natural homomorphism, mapping element from $K[x]$ to the
representative in $K'$. We shall prove that indeed $K$ is a subfield of $K'$.
For that, we shall prove that $\iota|_K \cong K$. Indeed, suppose $\gamma \in
K$. If $\iota(\gamma)=0$, then $\gamma \in \langle f(x) \rangle$. In other
words, $\gamma = f(x)h(x)$ for some $h \in K[x]$. Since $\deg f>0$, the only
possibility of this is when $\gamma=0$. Therefore, $\iota|_K$ is one-to-one.
Since it is obviously surjective, it is a bijection (and thus an isomorphism). 

To finish the proof, we construct $\alpha$. Take $\alpha := x + \langle f(x)
\rangle$ in $K'$. Then, $f(\alpha)=f(x+\langle f(x) \rangle) = f(x) = 0$ in
$K'$. Thus, we have proven the proposition. \qedwhite

Further, denote by $K(\alpha)$ the field $K'$ constructed as in the proof above.

\begin{proposition}\label{proposition:alphas}
    The elements $1,\alpha,\dots,\alpha^{n-1}$ form a basis for $K(\alpha)/K$
    with $n = \deg f$.
\end{proposition}

\textbf{Proof.} First we show linear independence. Suppose that we have a linear
combination that vanishes: $\sum_{j=0}^{n-1}c_j\alpha^j=0$. Define polynomial
$h(x) := \sum_{j=0}^{n-1}c_jx^j$. Then, we have $h(\alpha)=0$. Then, either $h$
is a multiple of $f$ or $0$. Since degree of $f$ is $n$, we conclude that $h$ is
a zero polynomial, so the set is linearly independent.

This set spans $K(\alpha)$. Indeed, consider any element $g(x)+\langle f(x)
\rangle \in K(\alpha)$. If $\deg g \geq n$, we can divide by $f(x)$, and the
remainder $r(x)$ will be equal to this element. In such case, $r(x) =
\sum_{j=0}^{n-1}r_jx^j$. Evaluate at $\alpha$:
$r(\alpha)=\sum_{j=0}^{n-1}r_j\alpha^j$. This is exactly the representation of
$g$ in $K(\alpha)$ using basis. \qedwhite

\subsection{Existence}\label{section:existence}

We prove the following key fact.
\begin{proposition}\label{theorem:product-irreducible} For integer $d$, we denote by
    $\Phi_d(x)$ the product of all monic irreducible polynomials of degree
    $d$ in $\mathbb{F}_{p}[x]$. Then, for every $n \geq 1$ the following holds:
    \begin{equation*}
        x^{p^n}-x = \prod_{d \mid n}\Phi_d(x)
    \end{equation*}
\end{proposition} 

\textbf{Proof.} Recall that $\mathbb{F}_p[x]$ is a UFD, and let $S$ be the set
of all monic irreducible polynomials with coefficients in $\mathbb{F}_p$. Let us
consider the factorization of $x^{p^n}-x$:
\begin{equation*}
    x^{p^n}-x = c\prod_{f \in S}f(x)^{e(f)}, \quad c \in \mathbb{F}_p^{\times}.
\end{equation*}

Note that $c=1$ since $x^{p^n}-x$ is monic. So all we have to prove is the
following:
\begin{itemize}
    \item[\textit{(i)}] If $f(x)$ divides $x^{p^n}-x$, then $f(x)^2$ does not
    divide $x^{p^n}-x$. This way, $e(f)$ is either $1$ or $0$.
    \item[\textit{(ii)}] Irreducible polynomial $f(x)$ divides $x^{p^n}-x$ if
    and only if $d = \deg f$ divides $n$. This way, $e(f)=0$ for $d \not\mid n$.
\end{itemize}

These two facts imply that indeed $x^{p^n}-x = \prod_{d \mid n}\Phi_d(x)$.

\textcolor{blue!70!black}{\textit{Part (i).}} Suppose $f(x) \mid (x^{p^n}-x)$. Then, $f(x)^2 \not\mid
(x^{p^n}-x)$. Indeed, suppose otherwise. Then, $f(x)^2g(x) = x^{p^n}-x$ for some
$g \in \mathbb{F}_p[x]$. Take the formal derivative of both sides:
\begin{equation*}
    2f(x)f'(x)g(x) + f(x)^2g'(x) = p^n x^{p^n-1} - 1 = -1 \quad \textcolor{gray!75!black}{\textit{(over $\mathbb{F}_p[x]$)}}
\end{equation*}
Thus, $f \mid 1$ which is not possible if $\deg f > 0$. Thus, $f(x)^2 \not\mid
(x^{p^n}-x)$.

\textcolor{blue!70!black}{\textit{Part (ii).}} So it remains to prove that if $f
\in S$, then one has $f(x) \mid (x^{p^n}-x) \iff d \mid n$. To prove this, we
construct $K := \mathbb{F}_p(\alpha)$ where $\alpha$ is a root of $f(x)$. By
Proposition~\ref{proposition:alphas}, $[K\!:\!\mathbb{F}_p]=d$, so $\#K = p^d$
and every element of $K$ is a root of polynomial $x^{p^d}-x$.

($\Rightarrow$) Assume that $x^{p^n}-x=f(x)g(x)$. Substitute $x=\alpha$. Then
one gets $\alpha^{p^n}=\alpha$. Consider any $\gamma :=
\sum_{j=0}^{d-1}b_j\alpha^j$ --- an element of $K$, then $\gamma^{p^n} =
\sum_{j=0}^{d-1}b_j(\alpha^{p^n})^j = \sum_{j=0}^{d-1}b_j\alpha^j = \gamma$
(here we use the fact that $\pi_{p^n}: x \mapsto x^{p^n}$ is a homomorphism).
Hence, elements of $K$ satisfy $x^{p^n}-x=0$. It follows that $(x^{p^d}-x) \mid
(x^{p^n}-x)$ and thus by Proposition~\ref{proposition:ell-m}, $d \mid n$.

($\Leftarrow$) Assume now that $d \mid n$. Since $\alpha^{p^d} = \alpha$ and $f(x)$ is the monic 
irreducible polynomial for $\alpha$, we have $f(x) \mid (x^{p^d}-x)$. Since $d \mid n$,
we have $(x^{p^d}-x) \mid (x^{p^n}-x)$ again by Proposition~\ref{proposition:ell-m}.
Thus $f(x) \mid (x^{p^n}-x)$. \qedwhite

Denote by $N_d$ the number of monic irreducible polynomials of degree $d$ in
$\mathbb{F}_p[x]$. Notice that $\deg \Phi_d = dN_d$. In such case, equating degrees 
of Proposition~\ref{theorem:product-irreducible}, we obtain
\begin{proposition}\label{proposition:nd} One has $p^n = \sum_{d \mid n}dN_d$.
    In particular, applying Möbius inversion formula, one obtains $N_n =
    n^{-1}\sum_{d \mid n}\mu(n/d)p^d$ where $\mu(n)$ is a Möbius function.
\end{proposition}

\begin{theorem}
    For each $n \geq 1$, there exists an irreducible polynomial over
    $\mathbb{F}_p$ of degree $n$.
\end{theorem}

\textbf{Proof.} From Proposition~\ref{proposition:nd}, the number of such
polynomials is $n^{-1}(p^n-\dots + p\mu(n))$. Consider the sum in 
the parenthesses:
\begin{equation*}
    S_n = p^n + R_n, \quad R_n := \sum_{j=1}^{n-1}b_jp^j, \quad b_j \in \{0,\pm 1\}.
\end{equation*}

We bound the expression for $R_n$:
\begin{equation*}
    |R_n| \leq \sum_{j=1}^{n-1}|b_j|p^j \leq \sum_{j=1}^{n-1}p^j \leq \frac{p^n-p}{p-1} < p^n
\end{equation*}

Thus, $S_n>0$, so $n^{-1}S_n>0$, thus we are done. \qedwhite

\subsection{Uniqueness}

\begin{lemma}\label{lemma:minimal}
    Let $p$ be a prime and $K$ be a finite field of order $p^n$. There exists an
    irreducible polynomial $f \in \mathbb{F}_p[x]$ such that $K \cong
    \mathbb{F}_p[x]/\langle f(x) \rangle$. 
\end{lemma}

\textbf{Proof.} Let $\alpha$ be a primitive element of $K$, i.e. a generator of
$K^{\times}$. Such an element exists by Theorem~\ref{theorem:fp-cyclic}. Since
$x^{p^n}-x=\prod_{d \mid n}\Phi_d(x)$, there is some irreducible polynomial $f
\in \mathbb{F}_p[x]$ such that $f(\alpha)=0$. But in such case, $\alpha^p,
\alpha^{p^2},\dots,\alpha^{p^{n-1}}$ are its distinct roots (they are distinct
because $\alpha$ has an order $p^n-1$). Indeed, $f(\alpha)^p = \sum_{j=0}^n
f_j^p\alpha^p = \sum_{j=0}^n f_j\alpha^p = f(\alpha^p)$. Thus, $f(\alpha^p)=0$.
By applying this repeatedly, we prove that the rest are also roots of $f$.
Therefore, $f$ is of degree $n$, thus $\mathbb{F}_p[x]/\langle f(x) \rangle$ is
a field with $p^n$ elements. Since both $\mathbb{F}_p[x]/\langle f(x) \rangle$
and $K$ are $n$-dimensional vector spaces over $\mathbb{F}_p$, they are
isomorphic. \qedwhite

\begin{theorem}
    Let $p$ be a prime. If $K_0$ and $K_1$ are both finite fields 
    of order $p^n$, then $K_0 \cong K_1$.
\end{theorem}

\textbf{Proof.} We know that $\mathbb{F}_p$ is a subfield of both $K_0$ and
$K_1$. From Lemma~\ref{lemma:minimal}, there exists an irreducible polynomial $f
\in \mathbb{F}_p[x]$ of degree $n$ such that $K_0 \cong \mathbb{F}_p[x]/\langle
f(x) \rangle$. But since all elements of $K_0$ are roots of $x^{p^n}-x$, we have
$f(x) \mid (x^{p^n}-x)$. But note that all elements in $K_1$ also satisfy
$x^{p^n}-x$, so there is some $\beta \in K_1$ such that $f(\beta)=0$. It follows
that $K_1 = \mathbb{F}_p(\beta)$ is a vector space of dimensionality $n$, so
$K_0 \cong K_1$. \qedwhite

% Define map $\psi: \mathbb{F}_p[x] \to K_1$ as $\sum_i c_ix^i \mapsto \sum_i
% c_i\beta^i$. Note that $\langle f \rangle \subseteq \text{ker}(\psi)$. But since
% $f$ is irreducible, $\langle f \rangle = \text{ker}(\psi)$. Thus, we have a
% monomorphism $\psi^*: \mathbb{F}_p[x] / \langle f(x) \rangle \hookrightarrow
% K_1$. But since $|\mathbb{F}_p[x] / \langle f(x) \rangle| = |K_1| = p^n$, it is
% a bijection. So in fact, $\psi^*$ is an isomorphism. Thus, $K_1 \cong
% \mathbb{F}_p[x]/\langle f(x) \rangle \cong K_0$, which is exactly what we had to
% prove. \qedwhite

\begin{example}
    Suppose $p \geq 5$. Then, there are at least $2$ quadratic non-residues 
    in $\mathbb{F}_p$: say, $a$ and $b$. Define two quadratic extensions:
    \begin{equation*}
        K_1 := \mathbb{F}_p[u] / \langle u^2-a\rangle, \quad K_2 := \mathbb{F}_p[v] / \langle v^2 - b \rangle.
    \end{equation*}
    One easily checks these are two finite fields of order $p^2$. Let us build
    the concrete isomorphism $\phi: K_1 \to K_2$. We need to understand what 
    $u$ gets mapped to. So let $\phi: u \mapsto \alpha + \beta v$. Then:
    \begin{equation*}
        \phi(u)^2 = (\alpha + \beta v)^2 = \alpha^2 + 2\alpha\beta v+ \beta^2 v^2 = (\alpha^2 + b\beta^2) + 2\alpha\beta v = a
    \end{equation*}

    Thus, $\alpha\beta=0$. If $\beta=0$, then $\alpha^2=a$, which has no
    solutions over $\mathbb{F}_p$. Thus, let $\alpha=0$, then $b\beta^2=a$, so
    $\beta=\sqrt{a/b}$. So the isomorphism is thus:
    \begin{equation*}
        \phi: c_0 + c_1u \mapsto c_0 + c_1\sqrt{a/b}\cdot v
    \end{equation*}
\end{example}

\subsection{Subfields \& Applications}

Important corollary which we have not yet mentioned explicitly is the following. 

\begin{lemma}\label{lemma:subfields} 
    Let $K$ be a finite field of dimensionality $n$
    over $\mathbb{F}_p$. The subfields of $K$ are in one-to-one correspondence
    with divisors $n$. In other words, if $d \mid n$, there exists a subfield $E
    \subseteq K$ of order $p^d$ and it is unique.
\end{lemma}

\textbf{Proof.} Suppose $E$ is a subfield of $K$ of dimensionality $d$ over
$\mathbb{F}_p$. We shall show that $d \mid n$.

Since $E^{\times}$ contains $p^d-1$ elements all satisfying $x^{p^d-1}-1$, we
have that $(x^{p^d-1}-1) \mid (x^{p^n-1}-1)$. By
Proposition~\ref{proposition:ell-m}, we have $(p^d-1) \mid (p^n-1)$. Applying
Proposition~\ref{proposition:ell-m} once again, $d \mid n$.

Now, from the theory above, this field exists and it is unique up to an
isomorphism.\qedwhite

Lemma~\ref{lemma:subfields} plays a central role in various practical
applications. For instance, in Cryptography, it is highly convenient and optimal
to work with low-bit prime fields: for instance, over the Mersenne prime
$p=2^{31}-1$. However, this field cannot be used in real-world applications,
since it is too small (breaking the protocol with the effective brute-force size
of $31$ bits is trivial for modern computers). There are two natural ways to
extend the number of elements: the first is to choose different $p$. However,
another practical solution is to construct extension $\mathbb{F}_{p^d}$ for
suitable $d$. For instance, for 31-bit Mersenne prime and $d=4$, one can
construct the 124-bit field extension $\mathbb{F}_{p^4}$, where the base field
arithmetic is done very effectively.

Yet, sometimes dealing with the finite field extensions is the necessity. For
instance, when implementing the bilinear pairing operation over elliptic curves,
by construction one has to perform arithmetic over $\mathbb{F}_{p^k}$ where $k$,
called an \textit{embedding degree}, is the minimal integer such that
$\mathbb{F}_{p^k}$ contains all $r$-th roots of unity (with $r$ being the order
of an elliptic curve). For suitable pairs $(p,r)$, the embedding degree is small
enough to be efficiently implemented.

Still, for now we simply stated importance of finite field extensions. However,
what is the use of Lemma~\ref{lemma:subfields}? Suppose it so happens that
$k=12$. In such case, what is the most efficient way to implement arithmetic
over $\mathbb{F}_{p^{12}}$? Obvious answer is to pick irreducible polynomial
$f(x) \in \mathbb{F}_{p}[x]$ of degree $12$ and conduct polynomial operations in
the base polynomial ring $\mathbb{F}_p[x]$ modulo $f(x)$. However, this sounds
highly suboptimal! The more efficient way of doing this is to exploit
Lemma~\ref{lemma:subfields}. We show this using a widely used prime field
from~\cite{bn254}.

\begin{example}[Barreto-Naehrig Field~\cite{bn254}]
    We consider the special case where the prime is parameterized as $p = 36x^4
    + 36x^3 + 24x^2 + 6x + 1$ with $x=\mathsf{0x44E992B44A6909F1}$.

    Turns out that for such prime $p$ and suitable elliptic curve over
    $\mathbb{F}_p$ (also parameterized by $x$), the embedding degree $k=12$.
    According to Lemma~\ref{lemma:subfields}, we have the following lattice for
    $\mathbb{F}_{p^{12}}$:

    \begin{center}
        \begin{tikzpicture}[
            node distance=10mm and 7.5mm,
            %every node/.style={draw, circle, minimum size=8mm, inner sep=0pt},
            >=stealth, thick
        ]

        % Nodes (from left to right)
        \node (p) {$\mathbb{F}_p$};
        \node (p2) [right=of p] {$\mathbb{F}_{p^2}$};
        \node (p3) [above=of p2] {$\mathbb{F}_{p^3}$};
        \node (p4) [right=of p2] {$\mathbb{F}_{p^4}$};
        \node (p6) [above=of p4] {$\mathbb{F}_{p^6}$};
        \node (p12) [right=of p6] {$\mathbb{F}_{p^{12}}$};

        % Edges
        \draw[->] (p) -- (p2);
        \draw[->] (p) -- (p3);
        \draw[->,red!80!purple] (p2) -- (p4);
        \draw[->,blue!80!black] (p2) -- (p6);
        \draw[->] (p3) -- (p6);
        \draw[->,red!80!purple] (p4) -- (p12);
        \draw[->,blue!80!black] (p6) -- (p12);

        \end{tikzpicture}
    \end{center}

    This way, instead of directly building extension $\mathbb{F}_p \subset
    \mathbb{F}_{p^{12}}$, we can go in three different ways:
    \begin{itemize}
        \item[\textit{(i)}] $\mathbb{F}_p \subset \mathbb{F}_{p^3} \subset \mathbb{F}_{p^6} \subset \mathbb{F}_{p^{12}}$.
        \item[\textit{(ii)}] $\mathbb{F}_p \subset \mathbb{F}_{p^2} \subset \mathbb{F}_{p^6} \subset \mathbb{F}_{p^{12}}$ (marked in \textcolor{blue!80!black}{blue}).
        \item[\textit{(iii)}] $\mathbb{F}_p \subset \mathbb{F}_{p^2} \subset
        \mathbb{F}_{p^4} \subset \mathbb{F}_{p^{12}}$ (marked in \textcolor{red!80!purple}{red}). 
    \end{itemize}

    Turns out (according, e.g., to~\cite{tower}) that the last two options are
    the most optimal: option \textcolor{red!80!purple}{\textit{(i)}} provides a
    better compression ratio while \textcolor{blue!80!black}{\textit{(ii)}} is
    slightly faster to implement. In particular, option
    \textcolor{blue!80!black}{\textit{(ii)}} is practically implemented using
    the following construction, where $\xi := 9+u \in \mathbb{F}_{p^2}$:
    \begin{align*}
        \mathbb{F}_{p^2} &= \mathbb{F}_p[u] / \langle u^2+1 \rangle, \\
        \mathbb{F}_{p^6} &= \mathbb{F}_{p^2}[v] / \langle v^3 - \xi \rangle, \\
        \mathbb{F}_{p^{12}} &= \mathbb{F}_{p^6}[w] / \langle w^2 - v \rangle,
    \end{align*}
\end{example}

\subsection*{References}
\footnotesize
\renewcommand{\refname}{\vspace{-1em}}
\bibliographystyle{alphaurl}
\bibliography{refs}

\end{document}